\section*{A physical stopping theorem for point-defect modulation hierarchies}

All Fourier transforms below use the uniform probability measure on a
Boolean cube, all logarithms are base two, and $0\log(1/0)=0$.

\paragraph{Definition.}
Let $f_0:\{-1,1\}^{n_0}\to\{-1,1\}$ be a linear threshold function.
Inductively, suppose $f_{\ell-1}$ is a Boolean function of $N_{\ell-1}$
variables.  Introduce a fresh block $Z_\ell\in\{-1,1\}^{m_\ell}$ and put
\[
 f_\ell(u,z)=f_{\ell-1}
 \bigl(u_1q_{m_\ell}(z),\ldots,
       u_{N_{\ell-1}}q_{m_\ell}(z)\bigr),
\]
where $q_m$ equals $-1$ at one point of the $m$-cube and $+1$ elsewhere.
Blocks and old variables are physically disjoint.

\paragraph{Theorem.}
For every depth $L$ and every positive block dimensions $m_1,\ldots,m_L$,
there is a physical restriction tree which queries $Z_L,Z_{L-1},\ldots,Z_1$
in that order and has the following properties.
\begin{enumerate}
\item Every leaf is obtained from $f_0$ by input sign flips and hence is an
LTF with the same squared Fourier spectrum as $f_0$.
\item If $P_\ell$ is the Fourier mass of $f_{\ell-1}$ on odd-cardinality
sets, then
\[
 H(f_L)=H(f_0)+\sum_{\ell=1}^L P_\ell H(q_{m_\ell}),\qquad
 I(f_L)=I(f_0)+\sum_{\ell=1}^L P_\ell I(q_{m_\ell}).
\]
The summands on the right are precisely the reach-weighted path entropy and
queried influence payments at level $\ell$; no variable or influence is used
twice.
\item For every $m\ge1$, $H(q_m)<4I(q_m)$.  Consequently, if
$H(f_0)\le C_{\rm LTF}I(f_0)$, then
\[
 H(f_L)\le \max\{4,C_{\rm LTF}\}I(f_L).
\]
Using the published all-LTF constant
$C_{\rm LTF}=\log_2(4+e)+2\log_2(e\pi)<8.937$, the hierarchy has the same
dimension- and depth-independent constant.
\end{enumerate}

\paragraph{Proof.}
For a Boolean $g$ on $N$ variables and Boolean $h$ on a fresh block, set
$F(u,z)=g(u_1h(z),\ldots,u_Nh(z))$.  Expanding $g$ gives
\[
 F(u,z)=\sum_{S\subseteq[N]}\widehat g(S)\chi_S(u)h(z)^{|S|}.
\]
Thus even $S$ retain the constant modulator and odd $S$ acquire $h$.
Orthogonality between the $u$-characters yields
\[
 H(F)=H(g)+P_{\rm odd}(g)H(h),\qquad
 I(F)=I(g)+P_{\rm odd}(g)I(h).
\]
The influence identity includes all old-coordinate influence in $I(g)$ and
all fresh-block influence in the second term.  Iteration proves the two
displayed telescopes.

After fixing $z$, $F(u,z)$ is $g$ with either all input signs retained or all
flipped.  Querying the newest block first therefore leaves another hierarchy
of depth one less, up to input sign flips.  Continuing reaches an input-sign
flip of $f_0$.  LTFs are closed under such flips (negate the corresponding
weights), proving terminal membership rather than merely assigning a label
to the leaves.

It remains to bound the path charge.  Write $N=2^m$ and $r=2/N$.  The empty
Fourier coefficient of $q_m$ is $1-r$, while each nonempty coefficient has
magnitude $r$.  Hence
\[
 I(q_m)=mr,
\]
and
\[
 H(q_m)=-(1-r)^2\log_2(1-r)^2
       +r(2-r)(2m-2).
\]
Using $-\log(1-r)\le r/(1-r)$,
\[
 -(1-r)^2\log_2(1-r)^2
 \le {2r(1-r)\over\ln2}<4r.
\]
Also
\[
 4mr-r(2-r)(2m-2)=4r+r^2(2m-2)>4r.
\]
Therefore $H(q_m)<4mr=4I(q_m)$.  Multiplying by $P_\ell\ge0$, summing, and
combining with the terminal LTF estimate proves the theorem.\qed

\paragraph{Scope.}
The theorem is all-dimensional and permits arbitrary depth and block sizes,
but it is not a universal stopping theorem.  Its decisive hypothesis is the
common modulator, which makes every physical branch an input-sign flip of the
same predecessor.  Bent splices and general unequal-square-spectrum branches
remain outside the result.  Full FEI remains open.
